\documentclass{article}
\usepackage[T1]{fontenc}
\usepackage{lmodern}
\usepackage{graphicx}
\usepackage{amsmath,amssymb}
\usepackage[a4paper,margin=2cm]{geometry}
\usepackage[absolute,overlay]{textpos}
\setlength{\TPHorizModule}{1mm}
\setlength{\TPVertModule}{1mm}
\usepackage[indonesian]{babel}
\usepackage{hyperref}
\setlength{\emergencystretch}{2em}
% unit: Halaman 49

\begin{document}

% === Halaman 49 (llm) ===

\textbf{Contoh 3:}

\begin{tabular}{lll}
Plainteks & : & \texttt{\color{red}crypto}isshortfor\texttt{\color{red}crypto}graphy} \\
Kunci & : & \texttt{abcdabcdabcdabcdabcdabcdabcdabcd} \\
Cipherteks & : & \texttt{\textbf{CSASTP}KVSIQUTGQU\textbf{CSASTP}IUAQJB} \\
\end{tabular}

\vspace{10mm}

\begin{itemize}
    \item Pada contoh ini, \texttt{crypto} dienkripsi menjadi kriptogram yang sama, yaitu \textbf{CSATP}.
    \item Tetapi kasus seperti ini tidak selalu demikian, misalnya pada contoh berikut ini....
\end{itemize}

\end{document}
